\section*{Capítulo \ref{ch:b2:conjugate-additions} --- Carbonilos conjugados: Michael y Robinson}

\begin{solution}{exo:b2:conjugate-additions:1}
\begin{itemize}
  \item \ce{CH3MgBr}: principalmente 1,2 (duro, irreversible), 1-metilciclohex-2-en-1-ol.
  \item \ce{(CH3)2CuLi}: 1,4, 3-metilciclohexanona.
  \item Tiofenol con base: 1,4 (el tiolato es blando), 3-(fenilsulfanil)ciclohexanona.
  \item \ce{LiAlH4}: 1,2, ciclohex-2-en-1-ol.
\end{itemize}
\end{solution}

\begin{solution}{exo:b2:conjugate-additions:2}
Dadores: pentano-2,4-diona, nitrometano, propanodioato de dietilo (\ce{C-H} ácido). Aceptores:
propenonitrilo, propenoato de metilo, but-3-en-2-ona (\ce{C=C} conjugado con un grupo aceptor).
\end{solution}

\begin{solution}{exo:b2:conjugate-additions:3}
\ce{C4H9Br + 2Li -> C4H9Li + LiBr}, y luego \ce{2C4H9Li + CuI -> (C4H9)2CuLi + LiI}, en éter a baja
temperatura bajo gas inerte.
\end{solution}

\begin{solution}{exo:b2:conjugate-additions:4}
(3-Oxobutil)propanodioato de dietilo, \ce{(EtOOC)2CHCH2CH2COCH3}. La hidrólisis da el ácido propanodioico
sustituido, que pierde \ce{CO2} al calentar (\cref{prop:b2:enolates-aldol:decarboxylation}):
ácido 5-oxohexanoico, \ce{HOOCCH2CH2CH2COCH3}.
\end{solution}

\begin{solution}{exo:b2:conjugate-additions:5}
Dimetilcuprato de litio en éter a baja temperatura, y luego tratamiento acuoso: \omterm{def:b2:conjugate-additions:addition-modes}{adición conjugada}.
El bromuro de metilmagnesio, más duro, se adiciona sobre todo en 1,2 y da 1-metilciclohex-2-en-1-ol como producto
principal.
\end{solution}

\begin{solution}{exo:b2:conjugate-additions:6}
La \ce{Et3N} desprotona un poco de \ce{EtSH} y da \ce{EtS-}; el tiolato se adiciona al carbono $\beta$, y el
\omterm{def:b2:enolates-aldol:enolate}{enolato} formado toma el protón de otra molécula de \ce{EtSH}, lo que regenera \ce{EtS-}. Producto:
4-(etilsulfanil)butan-2-ona. El azufre es grande y polarizable, y su par no enlazante es alto en energía: un \omterm{def:b2:frontier-orbitals:hard-soft}{nucleófilo
blando}, bajo \omterm{def:b2:frontier-orbitals:control}{control orbital}, que ataca donde el \omterm{def:b2:frontier-orbitals:frontier}{LUMO} es mayor; además, la adición es
reversible, lo que también favorece el aducto 1,4, más estable.
\end{solution}

\begin{solution}{exo:b2:conjugate-additions:7}
El cianuro ataca más deprisa el carbono carbonílico, más positivo: la cianhidrina es el producto cinético. Su
formación es reversible; a temperatura más alta, con tiempo, el cianuro se libera y vuelve a adicionarse en el
carbono $\beta$, dando el aducto que conserva el enlace $\pi$ \ce{C=O}, más fuerte: el producto termodinámico.
\end{solution}

\begin{solution}{exo:b2:conjugate-additions:8}
Con la 2-metilciclohexano-1,3-diona: la cetona de Wieland--Miescher (figura de la \omterm{def:b2:conjugate-additions:robinson}{anulación de Robinson}).
Con la ciclohexanona: 4,4a,5,6,7,8-hexahidronaftalen-2(3\textit{H})-ona, una enona bicíclica cuyo
\ce{C=C} termina en un carbono de fusión de los anillos, sin grupo metilo angular.
\end{solution}

\begin{solution}{exo:b2:conjugate-additions:9}
\ce{CH3COCH2CH2CH2COCH3}: C2 y C6 son los carbonilos, cinco carbonos de C2 a C6 contados ambos inclusive:
una 1,5-dicetona. Córtese C3--C4: dador, el \omterm{def:b2:enolates-aldol:enolate}{enolato} de la propanona en C3 (mejor, el 3-oxobutanoato de etilo, cuyo
éster se elimina después por hidrólisis y \omterm{def:b2:enolates-aldol:decarboxylation}{descarboxilación}); aceptor, la but-3-en-2-ona. Condiciones:
3-oxobutanoato de etilo, but-3-en-2-ona, etóxido de sodio catalítico en etanol; luego ácido acuoso y
calor.
\end{solution}

\begin{solution}{exo:b2:conjugate-additions:10}
La \omterm{def:b2:enolates-aldol:aldol}{aldólica} forma un enlace \ce{C-C} pero convierte un enlace $\pi$ \ce{C=O} en un enlace $\sigma$ \ce{C-O}; el
balance es pequeño y dos moléculas pasan a ser una, de modo que $K$ es modesta y la retroaldólica es fácil. La
\omterm{def:b2:conjugate-additions:michael}{adición de Michael} forma un enlace \ce{C-C} a costa de un enlace $\pi$ \ce{C=C}, más débil, y conserva ambos
\ce{C=O}: es claramente favorable. La reacción inversa exigiría que el \omterm{def:b2:enolates-aldol:enolate}{enolato} cetónico expulsara el anión propanodioato, estabilizado; el
equilibrio está tan desplazado hacia el aducto que, con base catalítica, el aducto persiste.
\end{solution}

\begin{solution}{exo:b2:conjugate-additions:11}
\ce{C(COOEt)2(CH2CH2COOCH3)2}. Tras la primera \omterm{def:b2:conjugate-additions:michael}{adición de Michael} el carbono central aún tiene un
hidrógeno entre dos ésteres, tan ácido como antes: la base lo arranca y sigue una segunda adición.
\end{solution}

\begin{solution}{exo:b2:conjugate-additions:12}
A través del \omterm{def:b2:enolates-aldol:enolate}{enolato} más sustituido (hacia el carbono que lleva el metilo): la octalona con el grupo metilo
en el carbono de fusión de los anillos, 4a-metil-4,4a,5,6,7,8-hexahidronaftalen-2(3\textit{H})-ona. A través
del \omterm{def:b2:enolates-aldol:enolate}{enolato} menos sustituido (del lado del \ce{CH2}): el grupo metilo acaba en un carbono del anillo contiguo a la
fusión. El primero está favorecido en condiciones de equilibrio (un alcóxido en su alcohol, o la vía de la \omterm{def:b2:amines:imine}{enamina}
con una \omterm{def:b2:amines:class}{amina secundaria} y un ácido), que dan el \omterm{def:b2:enolates-aldol:kinetic-enolate}{enolato termodinámico}, más sustituido
(\cref{prop:b2:enolates-aldol:enolate-control}).
\end{solution}

\begin{solution}{pb:b2:conjugate-additions:1}
\textbf{1.} Un anillo de ciclohexano con \ce{C=O} en C1 y C3 y un grupo metilo en C2; el hidrógeno de C2
es el más ácido.
\textbf{2.} Al arrancarlo se obtiene un anión cuya carga está deslocalizada sobre C2 y los dos oxígenos; en la
ciclohexanona solo la comparte un oxígeno, y la acidez es demasiado débil para medirse en agua.
\textbf{3.} Carga sobre C2; \ce{C1=C2} con \ce{O^-} sobre C1; \ce{C2=C3} con \ce{O^-} sobre C3.
\textbf{4.} Sí: con hidróxido, $K = 10^{14.0 - 5.3} = 10^{8.7}$ (el agua como ácido conjugado, a partir de
$K_e$); con etóxido, $10^{15.9-5.3} = 10^{10.6}$. Ambas desprotonaciones son completas.
% ledger: ke:25C, pka:ethanol
\textbf{5.} \ce{C1(OH)=C2(CH3)}: C2 aún lleva un hidrógeno, que es todo lo que necesita un \omterm{def:b2:enolates-aldol:tautomerism}{enol}; el grupo
metilo solo sustituye al segundo.
\textbf{6.} Blando: la carga está repartida sobre tres átomos y el anión está estabilizado. Ataca el
carbono $\beta$ (\cref{prop:b2:conjugate-additions:regio}).
\textbf{7.} El C2 del \omterm{def:b2:enolates-aldol:enolate}{enolato} se une al extremo \ce{CH2} de la but-3-en-2-ona; el \omterm{def:b2:enolates-aldol:enolate}{enolato} formado toma un
protón: 2-metil-2-(3-oxobutil)ciclohexano-1,3-diona.
\textbf{8.} Cada nuevo \omterm{def:b2:enolates-aldol:enolate}{enolato} toma un protón de una molécula de diona (o del disolvente) y así
regenera el \omterm{def:b2:enolates-aldol:enolate}{enolato} dador: la base no se consume.
\textbf{9.} El \omterm{def:b2:enolates-aldol:enolate}{enolato} es ambidentado; con un electrófilo carbonado blando, bajo \omterm{def:b2:frontier-orbitals:control}{control orbital}, reacciona
por el carbono, y el aducto \ce{C-C} conserva dos enlaces \ce{C=O}, la combinación más fuerte; un
aducto en \ce{O}, si se formara, revertiría.
\textbf{10.} Uno: C2, unido a C1, a C3, al metilo y a la cadena.
\textbf{11.} El carbono carbonílico de la cadena y el carbonilo del anillo C1 (o C3): \ce{C(=O)}--\ce{CH2}--\ce{CH2}--C2--C1,
cinco carbonos contados ambos extremos inclusive.
\textbf{12.} Para consumir toda la diona, el reactivo más valioso; solo ligero, porque la but-3-en-2-ona
es muy tóxica y polimeriza.
\textbf{13.} Los C4 y C6 del anillo (contiguos a C3 y C1), el \ce{CH2} de la cadena contiguo a su carbonilo, y el
\ce{CH3} terminal de la cadena. C2 no tiene ninguno.
\textbf{14.} El \omterm{def:b2:enolates-aldol:enolate}{enolato} del \ce{CH3} terminal ataca un carbonilo del anillo: el anillo cerrado cuenta ese
carbono, el carbono carbonílico de la cadena, dos \ce{CH2}, C2 y el carbono carbonílico atacado: seis átomos.
\textbf{15.} El \omterm{def:b2:enolates-aldol:enolate}{enolato} del \ce{CH2} interno de la cadena cerraría un anillo de cuatro eslabones, tensionado. Un \omterm{def:b2:enolates-aldol:enolate}{enolato}
del anillo (C4 o C6) que atacara el carbonilo de la cadena cerraría también un anillo de seis eslabones, pero con puente: ese \omterm{def:b2:enolates-aldol:aldol}{aldol}
no puede deshidratarse (el doble enlace quedaría en una cabeza de puente de un sistema bicíclico pequeño) y revierte,
mientras que el \omterm{def:b2:enolates-aldol:aldol}{aldol} fusionado se va retirando por deshidratación.
\textbf{16.} El \omterm{def:b2:enolates-aldol:aldol}{aldol} lleva el \ce{OH} sobre el antiguo carbono carbonílico del anillo, ahora en la fusión de los anillos; la pérdida
de agua con un hidrógeno del antiguo carbono del \ce{CH3} (ahora en $\alpha$ de la cetona de la cadena) da el
\ce{C=C} conjugado con esa cetona.
\textbf{17.} El nuevo doble enlace está conjugado con el carbonilo: una eliminación E1cB a través del
\omterm{def:b2:enolates-aldol:enolate}{enolato}, favorecida por el calentamiento (\cref{prop:b2:enolates-aldol:aldol-equilibrium}).
\textbf{18.} Una cetona saturada y una cetona $\alpha,\beta$-insaturada (una ciclohexenona).
\textbf{19.} El carbono de fusión de los anillos que lleva el grupo metilo: cuatro sustituyentes distintos.
\textbf{20.} No: la diona es aquiral, y sus dos carbonilos se atacan por igual (son imágenes especulares
uno del otro respecto a la cadena); con reactivos aquirales el producto es racémico. Los catalizadores de amina
quirales favorecen un carbonilo y dan un enantiómero en exceso (volumen del tercer año).
\textbf{21.} Diona \ce{C7H10O2}: \qty{126.155}{g/mol}; \ce{C4H6O}: \qty{70.091}{g/mol}; \ce{C11H14O2}:
\qty{178.231}{g/mol}.
\textbf{22.} \ce{C7H10O2 + C4H6O -> C11H14O2 + H2O}: C 11, H 16, O 3 a cada lado.
\textbf{23.} $10.0/126.155 = \qty{0.07927}{mol} = \qty{79.3}{mmol}$.
\textbf{24.} $0.07927 \times 178.231 \times 0.70 = \boldsymbol{\approx \qty{9.89}{g}}$ de
cetona de Wieland--Miescher.
\end{solution}
